\documentclass[10pt,a4paper]{amsart}
\usepackage{amsmath,amssymb,amsthm,graphicx}
\usepackage{verbatim}
\newenvironment{reference}{\comment}{\endcomment}
%\newenvironment{reference}{}{}
\newenvironment{slogan}{\comment}{\endcomment}
\newenvironment{history}{\comment}{\endcomment}

% For commutative diagrams we use Xy-pic
\usepackage[all]{xy}

% We use 2cell for 2-commutative diagrams.
\xyoption{2cell}
% Theorem environments.
%
\theoremstyle{plain}
\newtheorem{theorem}[subsection]{Theorem}
\newtheorem{proposition}[subsection]{Proposition}
\newtheorem{lemma}[subsection]{Lemma}

\theoremstyle{definition}
\newtheorem{definition}[subsection]{Definition}
\newtheorem{example}[subsection]{Example}
\newtheorem{exercise}[subsection]{Exercise}
\newtheorem{situation}[subsection]{Situation}

\theoremstyle{remark}
\newtheorem{remark}[subsection]{Remark}
\newtheorem{remarks}[subsection]{Remarks}

\numberwithin{equation}{subsection}

% Macros
%
\def\lim{\mathop{\mathrm{lim}}\nolimits}
\def\colim{\mathop{\mathrm{colim}}\nolimits}
\def\Spec{\mathop{\mathrm{Spec}}}
\def\Hom{\mathop{\mathrm{Hom}}\nolimits}
\def\Ext{\mathop{\mathrm{Ext}}\nolimits}
\def\SheafHom{\mathop{\mathcal{H}\!\mathit{om}}\nolimits}
\def\SheafExt{\mathop{\mathcal{E}\!\mathit{xt}}\nolimits}
\def\Sch{\mathit{Sch}}
\def\Mor{\mathop{\mathrm{Mor}}\nolimits}
\def\Ob{\mathop{\mathrm{Ob}}\nolimits}
\def\Sh{\mathop{\mathit{Sh}}\nolimits}
\def\NL{\mathop{N\!L}\nolimits}
\def\CH{\mathop{\mathrm{CH}}\nolimits}
\def\proetale{{pro\text{-}\acute{e}tale}}
\def\etale{{\acute{e}tale}}
\def\QCoh{\mathit{QCoh}}
\def\Ker{\mathop{\mathrm{Ker}}}
\def\Im{\mathop{\mathrm{Im}}}
\def\Coker{\mathop{\mathrm{Coker}}}
\def\Coim{\mathop{\mathrm{Coim}}}

% Boxtimes
%
\DeclareMathSymbol{\boxtimes}{\mathbin}{AMSa}{"02}

%
% Macros for moduli stacks/spaces
%
\def\QCohstack{\mathcal{QC}\!\mathit{oh}}
\def\Cohstack{\mathcal{C}\!\mathit{oh}}
\def\Spacesstack{\mathcal{S}\!\mathit{paces}}
\def\Quotfunctor{\mathrm{Quot}}
\def\Hilbfunctor{\mathrm{Hilb}}
\def\Curvesstack{\mathcal{C}\!\mathit{urves}}
\def\Polarizedstack{\mathcal{P}\!\mathit{olarized}}
\def\Complexesstack{\mathcal{C}\!\mathit{omplexes}}
% \Pic is the operator that assigns to X its picard group, usage \Pic(X)
% \Picardstack_{X/B} denotes the Picard stack of X over B
% \Picardfunctor_{X/B} denotes the Picard functor of X over B
\def\Pic{\mathop{\mathrm{Pic}}\nolimits}
\def\Picardstack{\mathcal{P}\!\mathit{ic}}
\def\Picardfunctor{\mathrm{Pic}}
\def\Deformationcategory{\mathcal{D}\!\mathit{ef}}
\usepackage{cleveref}
\crefname{originallemma}{original lemma}{original lemmas}
\Crefname{originallemma}{Original lemma}{Original lemmas}
\makeatletter
\newlabel{categories-lemma-what-gets-inverted}{{05Q9}{source}{}{sourcepointer-1}{}}
\newlabel{categories-lemma-what-gets-inverted@cref}{{[originallemma][0][]05Q9}{[1][0][]source}{}{}{}}
\makeatother
\title{The objects killed by the localization functor at a multiplicative system}
\author{The Stacks Project Authors}
\date{stacks-project:homology.tex:local-snapshot-20261012; GNU FDL 1.2 or later}
\begin{document}
\maketitle
\noindent\textit{Editorial note.} All delivered mathematical text is the author's own, copied from the pinned Stacks Project chapter file homology.tex; no statement or proof step was rewritten, reordered or removed.\par\medskip
\noindent\textit{Editorial note.} The proof is the author's complete argument for all three equivalences, each of which is written out with the morphism it applies the hypotheses to; no step is left to the reader and the closing sentence is a conclusion.\par\medskip
\noindent\textit{Editorial note.} Every labelled reference occurring inside the counted spans is carried in assembly.external\_references as a non-load-bearing navigation pointer (same-chapter ones at their definition site, cross-chapter ones at the invocation inside the counted span), so no load-bearing mathematics is replaced by a stub. No bibliography citation occurs inside the counted spans.\par\medskip
\noindent\textit{Editorial note.} The same-name check for this card scanned every staged Stacks chapter of this repository -- all sources/stacks-*/pinned-original-source.tex files, the card's own chapter excluded -- and the labels found are recorded in dedup.same\_name\_cross\_chapter on the card.\par\medskip
\noindent\textit{Editorial note.} The preamble is the chapter's own native preamble excerpt already used for this source: the comment-type reference, slogan and history environments (native preamble lines 8-13), the Xy-pic setup (native preamble lines 14-18) and the native theorem declarations with the mathematical macros (native preamble lines 154-217, the actual last line of the file).\par\medskip
\noindent\textit{Allowed standard prerequisites (editorial).}\begin{itemize}
\item Localization at a multiplicative system. Same chapter and Categories, the construction of the localization and the universal property of the localization functor. Supplies the functor Q and the fact that it carries the morphisms of the multiplicative system to isomorphisms.
\item Which morphisms the localization inverts. Categories, Lemma categories-lemma-what-gets-inverted. Supplies the criterion that a morphism becomes an isomorphism in the localization exactly when it can be completed to an element of the multiplicative system, which produces condition (3) from condition (1).
\end{itemize}\medskip
\ifdefined\hypertarget\hypertarget{sourcepointer-1}{}\fi
\noindent\textit{Source reference.} Original lemma 05Q9: sources/stacks-categories/pinned-original-source.tex, source lines 5013--5013. Navigation pointer invoked at proof line 2122; the counted span contains the pointer only and the cited result is not reproduced inside it, so the pointer is registered as a non-load-bearing navigation stub in the established form for Stacks material. No bibliography citation is involved.\par\medskip
\begin{lemma}
\label{lemma-kernel-localization}
Let $\mathcal{C}$ be an additive category. Let $S$ be a multiplicative
system. Let $X$ be an object
of $\mathcal{C}$. The following are equivalent
\begin{enumerate}
\item $Q(X) = 0$ in $S^{-1}\mathcal{C}$,
\item there exists $Y \in \Ob(\mathcal{C})$ such that
$0 : X \to Y$ is an element of $S$, and
\item there exists $Z \in \Ob(\mathcal{C})$ such that
$0 : Z \to X$ is an element of $S$.
\end{enumerate}
\end{lemma}

\begin{proof}
If (2) holds we see that $0 = Q(0) : Q(X) \to Q(Y)$ is an isomorphism.
In the additive category $S^{-1}\mathcal{C}$ this implies that $Q(X) = 0$.
Hence (2) $\Rightarrow$ (1). Similarly, (3) $\Rightarrow$ (1).
Suppose that $Q(X) = 0$. This implies that the morphism
$f : 0 \to X$ is transformed into an isomorphism in $S^{-1}\mathcal{C}$.
Hence by
Categories, Lemma \ref{categories-lemma-what-gets-inverted}
there exists a morphism $g : Z \to 0$ such that $fg \in S$. This proves
(1) $\Rightarrow$ (3). Similarly, (1) $\Rightarrow$ (2).
\end{proof}

\end{document}
